\chapter{Preface}
\label{ch:preface}

\newthought{This book documents \Jtwo}, the distributed runtime of the \Jarvis\
framework: a Python system for orchestrating likelihood-based parameter scans in
high-energy physics. \Jtwo\ drives many long-lived worker processes through a Redis task
queue, runs external calculators and in-process operators for every parameter point, and
persists everything through an asynchronous archiver --- while you describe the whole scan
in a single \textsc{yaml} \emph{task card}.

The manual is written to be read in two ways. Read Part~\ref{part:getting-started} front to
back once: it installs the software, runs a first scan, and builds the vocabulary the rest
of the book leans on. After that, treat the book as a reference --- every \textsc{yaml} key,
every \cli{Jarvis} subcommand, and every output file has a home chapter, and the
appendices provide flat lookup tables that point back into the prose.

\section{Who this book is for}

You should be comfortable with a Unix shell, a text editor, and elementary Python. No prior
exposure to any version of \Jarvis\ is assumed: everything the runtime does is explained
here, from scratch. Familiarity with the physics tooling you plan to orchestrate
(\textsc{spheno}, \textsc{madgraph}, \textsc{rivet}, private codes, \ldots) helps for
Part~\ref{part:backends}, but the built-in examples run without any of them.

\section{How the book is organised}

\begin{description}[style=nextline, leftmargin=1.2em]
  \item[Part~\ref{part:getting-started}: Getting Started] What \Jtwo\ is, installation
        (including the Redis service it depends on), a ten-minute quick start, and the core
        concepts: projects, tokens, u-space, and the Sample lifecycle.
  \item[Part~\ref{part:task-card}: The Task Card] A field-by-field reference for the
        \textsc{yaml} task card: how it is loaded, the shared expression language every
        \yamlkey{expression:} field speaks, the \yamlkey{Sampling} block, variables and
        distributions, the likelihood, and the runtime settings under
        \yamlkey{EnvReqs.V2}.
  \item[Part~\ref{part:backends}: Workflow Backends] How a parameter point becomes
        observables: the command system that turns a card into processes, the layered
        workflow model, external \yamlkey{Calculators} together with the file formats
        they exchange through \Portal, in-process \yamlkey{Operas}, and the
        \yamlkey{LibDeps} tool registry.
  \item[Part~\ref{part:samplers}: Samplers] The five sampling methods --- \sampler{Random},
        \sampler{Grid}, \sampler{CSV}, \sampler{Bridson}, and the feedback-driven
        \sampler{AdaptiveBridson} --- plus the \cli{check} smoke path.
  \item[Part~\ref{part:running}: Running, Monitoring \& Outputs] The \cli{Jarvis} command
        line, project tools and the official example catalog, checkpoint/resume, the output
        tree, live monitoring, the run summary, and logging.
  \item[Part~\ref{part:advanced}: Advanced \& Performance] The distributed runtime in
        depth (Redis schema, Workers, the Archiver), performance tuning, plotting with
        \JPlot, extending the expression language with \Operas, and reproducibility.
\end{description}

\section{Conventions}

\begin{itemize}
  \item \yamlkey{Sampling.Method} --- a \textsc{yaml} key or block, written as a dotted
        path from the top of the task card.
  \item A command you type in a shell appears in a titled box, exactly as you would
        type it:
\begin{terminal}
\begin{jcmd}
Jarvis run task.yaml
\end{jcmd}
\end{terminal}
  \item A short command mentioned mid-sentence uses \cli{Jarvis -{}-help} instead.
  \item \token{@SampleID}, \marker{\&J/} --- runtime tokens and path markers, coloured
        like a \textsc{yaml} value (Chapter~\ref{ch:core-concepts}; full table in
        Chapter~\ref{ch:command-system}).
  \item \sampler{Bridson} --- a sampler method name.
  \item \req\ / \opt\ --- whether omitting a key raises an error or silently applies a
        default.
\end{itemize}

Three kinds of boxes interrupt the text:

\begin{jnote}
A \textbf{note} adds background or records a contract you can rely on.
\end{jnote}

\begin{jtip}
A \textbf{tip} is advice from production use --- defaults to keep, habits that save time.
\end{jtip}

\begin{jwarn}
A \textbf{warning} marks behavior that can cost you data, hours of compute, or an evening
of debugging.
\end{jwarn}

\section{Versions and errata}

This manual is a living document revised in step with the software; this build documents
\Jarvis\ \Jversion. The authoritative, machine-checked inventory of every \textsc{yaml} key
lives with the source; where the book and the code disagree, the code wins --- and we would
appreciate an issue on the project tracker so the book can be fixed.
